% --- Archon physics formalization source begin ---
% archon:physics
% archon:covers PhyXMiniProblems/problem_phyx_mini_0530.lean
% archon:source-report reports/phyx_mini/problem_phyx_mini_0530.source.json

\chapter{Physics problem phyx\_mini\_0530}
\label{ch:PhyXMiniProblems_problem_phyx_mini_0530}

\paragraph{Problem source.}
\#\# Physical scenario

A photon having wavelength \textbackslash{}( \textbackslash{}lambda \textbackslash{}) scatters off a free electron at A as shown in figure, producing a second photon having wavelength \textbackslash{}( \textbackslash{}lambda' \textbackslash{}). This photon then scatters off another free electron at B, producing a third photon having wavelength \textbackslash{}( \textbackslash{}lambda'' \textbackslash{}) and moving in a direction directly opposite the original photon as shown in the figure.

\#\# Question

Determine the value of \textbackslash{}( \textbackslash{}Delta \textbackslash{}lambda = \textbackslash{}lambda'' - \textbackslash{}lambda \textbackslash{}).

\#\# Answer choices

- A: A: \textbackslash{}( 2.15 \textbackslash{}times 10\textasciicircum{}\{-12\} \textbackslash{}text\{ m\} \textbackslash{})

- B: B: \textbackslash{}( 3.66 \textbackslash{}times 10\textasciicircum{}\{-12\} \textbackslash{}text\{ m\} \textbackslash{})

- C: C: \textbackslash{}( 4.72 \textbackslash{}times 10\textasciicircum{}\{-12\} \textbackslash{}text\{ m\} \textbackslash{})

- D: D: \textbackslash{}( 4.85 \textbackslash{}times 10\textasciicircum{}\{-12\} \textbackslash{}text\{ m\} \textbackslash{})

\#\# Figure caption (auxiliary; use the image as primary evidence)

The image is a diagram illustrating a scattering event, typically related to quantum or particle physics. 

- **Waves and Particles:**
  - Red wavy arrows represent photons or electromagnetic waves labeled as \textbackslash{}( \textbackslash{}lambda \textbackslash{}), \textbackslash{}( \textbackslash{}lambda' \textbackslash{}), and \textbackslash{}( \textbackslash{}lambda'' \textbackslash{}).
  - Blue spheres labeled "Electron 1" and "Electron 2" represent electrons.

- **Interaction Points:**
  - Point \textbackslash{}( A \textbackslash{}): A photon with wavelength \textbackslash{}( \textbackslash{}lambda \textbackslash{}) interacts here, leading to an emitted electron (Electron 1) and a photon with modified wavelength \textbackslash{}( \textbackslash{}lambda' \textbackslash{}).
  - Point \textbackslash{}( B \textbackslash{}): The photon with wavelength \textbackslash{}( \textbackslash{}lambda' \textbackslash{}) undergoes further interaction, resulting in another emitted electron (Electron 2) and another wavelength \textbackslash{}( \textbackslash{}lambda'' \textbackslash{}).

- **Angles:**
  - Angle \textbackslash{}( \textbackslash{}alpha \textbackslash{}): Between the initial photon direction and the direction of Electron 1.
  - Angle \textbackslash{}( \textbackslash{}theta \textbackslash{}): Between the initial photon and the line connecting point \textbackslash{}( A \textbackslash{}).
  - Angle \textbackslash{}( \textbackslash{}beta \textbackslash{}): Between the path of photon \textbackslash{}( \textbackslash{}lambda'' \textbackslash{}) and the path of Electron 2.

- **Dashed Lines:**
  - These lines indicate angles and directions related to the electrons and photon paths, helping illustrate the angular relationships.

The diagram effectively describes the sequence and geometrical configuration of scattering events involving photons and electrons.

\#\# Dataset metadata

Quantum Phenomena; Physical Model Grounding Reasoning, Spatial Relation Reasoning

\paragraph{Recorded answer/context.}
D: D: \textbackslash{}( 4.85 \textbackslash{}times 10\textasciicircum{}\{-12\} \textbackslash{}text\{ m\} \textbackslash{})

\paragraph{Figure/image path.}
/root/proposal\_for\_physic/science-mango/phyx\_mini\_run/phyx\_data/test\_image/530.png

\paragraph{Formalization target.}
create a compiling Lean file with sorry bodies at `PhyXMiniProblems/problem\_phyx\_mini\_0530.lean`. The Lean declarations must preserve the physical quantities, dimensions or dimensional roles, figure labels, governing-law hypotheses, and final relation expressed by this problem.
Use Mathlib/Physlib names found through LeanExplore where available. If a physics API is missing, introduce faithful local abstractions rather than scalar placeholder aliases.

\begin{theorem}[Physics formalization target]
\label{thm:physics:phyx_mini_0530:target}
\lean{PhyXMiniProblems.ProblemPhyXMini0530.problem_phyx_mini_0530}
\uses{def:physics:phyx-mini-0530:phyxminiproblems-problemphyxmini0530-sequentialcomptonsetup, def:physics:phyx-mini-0530:phyxminiproblems-problemphyxmini0530-matchessequentialscatteringscenario, def:physics:phyx-mini-0530:phyxminiproblems-problemphyxmini0530-matchessequentialcomptonfigure, def:physics:phyx-mini-0530:phyxminiproblems-problemphyxmini0530-hasphysicalsequentialcomptonparameters, def:physics:phyx-mini-0530:phyxminiproblems-problemphyxmini0530-matcheselectroncomptonwavelengthreadout, def:physics:phyx-mini-0530:phyxminiproblems-problemphyxmini0530-satisfiestwostagecomptonscatteringlaws, def:physics:phyx-mini-0530:phyxminiproblems-problemphyxmini0530-wavelengthshiftinpicometers, def:physics:phyx-mini-0530:phyxminiproblems-problemphyxmini0530-answerchoice, def:physics:phyx-mini-0530:phyxminiproblems-problemphyxmini0530-matchesanswerchoice}
The assigned autoformalize agent should translate this physics problem into Lean declarations in the covered file, with theorem and lemma proof bodies written as `by sorry`.
\end{theorem}
\begin{proof}
This is an autoformalization task, not a proof task. Produce statements that can later be proved by the physics prover without weakening the physical model.
\end{proof}

% --- Archon declaration topology begin ---
\section{Declaration topology}
\begin{definition}[Length Quantity]
  \label{def:physics:phyx-mini-0530:phyxminiproblems-problemphyxmini0530-lengthquantity}
  \lean{PhyXMiniProblems.ProblemPhyXMini0530.LengthQuantity}
  A nonnegative, unit-independent physical wavelength or other length
\end{definition}
\begin{proof}
  This is the declared model definition of Length Quantity.
\end{proof}
\begin{definition}[Direction2 D]
  \label{def:physics:phyx-mini-0530:phyxminiproblems-problemphyxmini0530-direction2d}
  \lean{PhyXMiniProblems.ProblemPhyXMini0530.Direction2D}
  A direction vector in the plane of the scattering diagram
\end{definition}
\begin{proof}
  This is the declared model definition of Direction2 D.
\end{proof}
\begin{definition}[length Readout]
  \label{def:physics:phyx-mini-0530:phyxminiproblems-problemphyxmini0530-lengthreadout}
  \lean{PhyXMiniProblems.ProblemPhyXMini0530.lengthReadout}
  \uses{def:physics:phyx-mini-0530:phyxminiproblems-problemphyxmini0530-lengthquantity}
  Read a physical length in a selected unit
\end{definition}
\begin{proof}
  This is the declared model definition of length Readout.
\end{proof}
\begin{definition}[length In Picometers]
  \label{def:physics:phyx-mini-0530:phyxminiproblems-problemphyxmini0530-lengthinpicometers}
  \lean{PhyXMiniProblems.ProblemPhyXMini0530.lengthInPicometers}
  \uses{def:physics:phyx-mini-0530:phyxminiproblems-problemphyxmini0530-lengthquantity, def:physics:phyx-mini-0530:phyxminiproblems-problemphyxmini0530-lengthreadout}
  Read a physical wavelength in picometers (10⁻¹² m)
\end{definition}
\begin{proof}
  This is the declared model definition of length In Picometers.
\end{proof}
\begin{definition}[direction Angle]
  \label{def:physics:phyx-mini-0530:phyxminiproblems-problemphyxmini0530-directionangle}
  \lean{PhyXMiniProblems.ProblemPhyXMini0530.directionAngle}
  \uses{def:physics:phyx-mini-0530:phyxminiproblems-problemphyxmini0530-direction2d}
  The undirected angle, in radians, between two propagation directions
\end{definition}
\begin{proof}
  This is the declared model definition of direction Angle.
\end{proof}
\begin{definition}[Interaction Point]
  \label{def:physics:phyx-mini-0530:phyxminiproblems-problemphyxmini0530-interactionpoint}
  \lean{PhyXMiniProblems.ProblemPhyXMini0530.InteractionPoint}
  The two interaction points labeled in the primary figure
\end{definition}
\begin{proof}
  This is the declared model definition of Interaction Point.
\end{proof}
\begin{definition}[Photon Label]
  \label{def:physics:phyx-mini-0530:phyxminiproblems-problemphyxmini0530-photonlabel}
  \lean{PhyXMiniProblems.ProblemPhyXMini0530.PhotonLabel}
  The three photon paths labeled by their wavelengths in the figure
\end{definition}
\begin{proof}
  This is the declared model definition of Photon Label.
\end{proof}
\begin{definition}[Electron Label]
  \label{def:physics:phyx-mini-0530:phyxminiproblems-problemphyxmini0530-electronlabel}
  \lean{PhyXMiniProblems.ProblemPhyXMini0530.ElectronLabel}
  The two recoil electrons explicitly labeled in the figure
\end{definition}
\begin{proof}
  This is the declared model definition of Electron Label.
\end{proof}
\begin{definition}[Figure Angle Label]
  \label{def:physics:phyx-mini-0530:phyxminiproblems-problemphyxmini0530-figureanglelabel}
  \lean{PhyXMiniProblems.ProblemPhyXMini0530.FigureAngleLabel}
  The three angle annotations printed in the primary figure
\end{definition}
\begin{proof}
  This is the declared model definition of Figure Angle Label.
\end{proof}
\begin{definition}[Photon State]
  \label{def:physics:phyx-mini-0530:phyxminiproblems-problemphyxmini0530-photonstate}
  \lean{PhyXMiniProblems.ProblemPhyXMini0530.PhotonState}
  \uses{def:physics:phyx-mini-0530:phyxminiproblems-problemphyxmini0530-lengthquantity, def:physics:phyx-mini-0530:phyxminiproblems-problemphyxmini0530-direction2d}
  A photon state retains both its physical wavelength and propagation direction
\end{definition}
\begin{proof}
  This is the declared model definition of Photon State.
\end{proof}
\begin{definition}[Sequential Compton Figure]
  \label{def:physics:phyx-mini-0530:phyxminiproblems-problemphyxmini0530-sequentialcomptonfigure}
  \lean{PhyXMiniProblems.ProblemPhyXMini0530.SequentialComptonFigure}
  \uses{def:physics:phyx-mini-0530:phyxminiproblems-problemphyxmini0530-interactionpoint, def:physics:phyx-mini-0530:phyxminiproblems-problemphyxmini0530-photonlabel, def:physics:phyx-mini-0530:phyxminiproblems-problemphyxmini0530-electronlabel, def:physics:phyx-mini-0530:phyxminiproblems-problemphyxmini0530-figureanglelabel}
  Qualitative information visibly present in the supplied diagram. The incidence maps say that λ enters A, λ' leaves A and enters B, and λ'' leaves B; they assign no numerical wavelength
\end{definition}
\begin{proof}
  This is the declared model definition of Sequential Compton Figure.
\end{proof}
\begin{definition}[Sequential Compton Setup]
  \label{def:physics:phyx-mini-0530:phyxminiproblems-problemphyxmini0530-sequentialcomptonsetup}
  \lean{PhyXMiniProblems.ProblemPhyXMini0530.SequentialComptonSetup}
  \uses{def:physics:phyx-mini-0530:phyxminiproblems-problemphyxmini0530-lengthquantity, def:physics:phyx-mini-0530:phyxminiproblems-problemphyxmini0530-direction2d, def:physics:phyx-mini-0530:phyxminiproblems-problemphyxmini0530-photonlabel, def:physics:phyx-mini-0530:phyxminiproblems-problemphyxmini0530-electronlabel, def:physics:phyx-mini-0530:phyxminiproblems-problemphyxmini0530-figureanglelabel, def:physics:phyx-mini-0530:phyxminiproblems-problemphyxmini0530-photonstate, def:physics:phyx-mini-0530:phyxminiproblems-problemphyxmini0530-sequentialcomptonfigure}
  Independent physical quantities for the two-stage event. In particular, the three photon wavelengths are independent fields rather than definitions from the requested net shift or from an answer choice
\end{definition}
\begin{proof}
  This is the declared model definition of Sequential Compton Setup.
\end{proof}
\begin{definition}[Matches Sequential Scattering Scenario]
  \label{def:physics:phyx-mini-0530:phyxminiproblems-problemphyxmini0530-matchessequentialscatteringscenario}
  \lean{PhyXMiniProblems.ProblemPhyXMini0530.MatchesSequentialScatteringScenario}
  \uses{def:physics:phyx-mini-0530:phyxminiproblems-problemphyxmini0530-sequentialcomptonsetup}
  The photon and electron roles at the two pictured interaction points
\end{definition}
\begin{proof}
  This is the declared model definition of Matches Sequential Scattering Scenario.
\end{proof}
\begin{definition}[Matches Sequential Compton Figure]
  \label{def:physics:phyx-mini-0530:phyxminiproblems-problemphyxmini0530-matchessequentialcomptonfigure}
  \lean{PhyXMiniProblems.ProblemPhyXMini0530.MatchesSequentialComptonFigure}
  \uses{def:physics:phyx-mini-0530:phyxminiproblems-problemphyxmini0530-directionangle, def:physics:phyx-mini-0530:phyxminiproblems-problemphyxmini0530-sequentialcomptonsetup}
  Facts read from the primary image. The final photon is antiparallel to the initial photon. The angle θ is between the initial and intermediate photon paths, while α is the first recoil angle. In the image, β lies between the continued intermediate path and Electron 2's recoil path
\end{definition}
\begin{proof}
  This is the declared model definition of Matches Sequential Compton Figure.
\end{proof}
\begin{definition}[Has Physical Sequential Compton Parameters]
  \label{def:physics:phyx-mini-0530:phyxminiproblems-problemphyxmini0530-hasphysicalsequentialcomptonparameters}
  \lean{PhyXMiniProblems.ProblemPhyXMini0530.HasPhysicalSequentialComptonParameters}
  \uses{def:physics:phyx-mini-0530:phyxminiproblems-problemphyxmini0530-lengthreadout, def:physics:phyx-mini-0530:phyxminiproblems-problemphyxmini0530-sequentialcomptonsetup}
  Positivity and nondegeneracy conditions ensure every photon has a genuine wavelength and direction and both electron recoil directions are meaningful
\end{definition}
\begin{proof}
  This is the declared model definition of Has Physical Sequential Compton Parameters.
\end{proof}
\begin{definition}[Matches Electron Compton Wavelength Readout]
  \label{def:physics:phyx-mini-0530:phyxminiproblems-problemphyxmini0530-matcheselectroncomptonwavelengthreadout}
  \lean{PhyXMiniProblems.ProblemPhyXMini0530.MatchesElectronComptonWavelengthReadout}
  \uses{def:physics:phyx-mini-0530:phyxminiproblems-problemphyxmini0530-lengthinpicometers, def:physics:phyx-mini-0530:phyxminiproblems-problemphyxmini0530-sequentialcomptonsetup}
  The standard electron Compton wavelength, calibrated in picometers. The value 2.42631 pm is independent physical reference data, not the requested two-scattering shift or an answer choice
\end{definition}
\begin{proof}
  This is the declared model definition of Matches Electron Compton Wavelength Readout.
\end{proof}
\begin{definition}[Satisfies Two Stage Compton Scattering Laws]
  \label{def:physics:phyx-mini-0530:phyxminiproblems-problemphyxmini0530-satisfiestwostagecomptonscatteringlaws}
  \lean{PhyXMiniProblems.ProblemPhyXMini0530.SatisfiesTwoStageComptonScatteringLaws}
  \uses{def:physics:phyx-mini-0530:phyxminiproblems-problemphyxmini0530-lengthreadout, def:physics:phyx-mini-0530:phyxminiproblems-problemphyxmini0530-directionangle, def:physics:phyx-mini-0530:phyxminiproblems-problemphyxmini0530-sequentialcomptonsetup}
  The governing Compton law at each interaction. For an initially stationary free electron, the outgoing photon wavelength exceeds the incoming one by λ\_C (1 - cos φ), where φ is the angle between the photon paths. These are two local laws; neither field states the requested total shift
\end{definition}
\begin{proof}
  This is the declared model definition of Satisfies Two Stage Compton Scattering Laws.
\end{proof}
\begin{lemma}[photon Scattering Angles Are Supplementary]
  \label{lem:physics:phyx-mini-0530:phyxminiproblems-problemphyxmini0530-photonscatteringanglesaresupplementary}
  \lean{PhyXMiniProblems.ProblemPhyXMini0530.photonScatteringAnglesAreSupplementary}
  \uses{def:physics:phyx-mini-0530:phyxminiproblems-problemphyxmini0530-directionangle, def:physics:phyx-mini-0530:phyxminiproblems-problemphyxmini0530-sequentialcomptonsetup, def:physics:phyx-mini-0530:phyxminiproblems-problemphyxmini0530-matchessequentialcomptonfigure, def:physics:phyx-mini-0530:phyxminiproblems-problemphyxmini0530-hasphysicalsequentialcomptonparameters}
  Because the last photon is opposite the first, the two undirected photon scattering angles are supplementary. This is a consequence of the figure, not an answer-value premise
\end{lemma}
\begin{proof}
  The conclusion follows from the governing relations and figure readouts named in the statement of photon Scattering Angles Are Supplementary.
\end{proof}
\begin{definition}[wavelength Shift In Picometers]
  \label{def:physics:phyx-mini-0530:phyxminiproblems-problemphyxmini0530-wavelengthshiftinpicometers}
  \lean{PhyXMiniProblems.ProblemPhyXMini0530.wavelengthShiftInPicometers}
  \uses{def:physics:phyx-mini-0530:phyxminiproblems-problemphyxmini0530-lengthinpicometers, def:physics:phyx-mini-0530:phyxminiproblems-problemphyxmini0530-sequentialcomptonsetup}
  The requested scalar Δλ = λ'' - λ, read in picometers
\end{definition}
\begin{proof}
  This is the declared model definition of wavelength Shift In Picometers.
\end{proof}
\begin{lemma}[total Shift Equals Twice Electron Compton Wavelength]
  \label{lem:physics:phyx-mini-0530:phyxminiproblems-problemphyxmini0530-totalshiftequalstwiceelectroncomptonwavelength}
  \lean{PhyXMiniProblems.ProblemPhyXMini0530.totalShiftEqualsTwiceElectronComptonWavelength}
  \uses{def:physics:phyx-mini-0530:phyxminiproblems-problemphyxmini0530-lengthinpicometers, def:physics:phyx-mini-0530:phyxminiproblems-problemphyxmini0530-sequentialcomptonsetup, def:physics:phyx-mini-0530:phyxminiproblems-problemphyxmini0530-matchessequentialcomptonfigure, def:physics:phyx-mini-0530:phyxminiproblems-problemphyxmini0530-hasphysicalsequentialcomptonparameters, def:physics:phyx-mini-0530:phyxminiproblems-problemphyxmini0530-satisfiestwostagecomptonscatteringlaws, def:physics:phyx-mini-0530:phyxminiproblems-problemphyxmini0530-wavelengthshiftinpicometers}
  The local Compton laws and supplementary-angle geometry make the net shift exactly twice the electron Compton wavelength. No calibration or answer choice is used in this structural relation
\end{lemma}
\begin{proof}
  The conclusion follows from the governing relations and figure readouts named in the statement of total Shift Equals Twice Electron Compton Wavelength.
\end{proof}
\begin{definition}[Answer Choice]
  \label{def:physics:phyx-mini-0530:phyxminiproblems-problemphyxmini0530-answerchoice}
  \lean{PhyXMiniProblems.ProblemPhyXMini0530.AnswerChoice}
  Labels of the four wavelength-shift choices printed in the problem
\end{definition}
\begin{proof}
  This is the declared model definition of Answer Choice.
\end{proof}
\begin{definition}[displayed Shift In Picometers]
  \label{def:physics:phyx-mini-0530:phyxminiproblems-problemphyxmini0530-displayedshiftinpicometers}
  \lean{PhyXMiniProblems.ProblemPhyXMini0530.displayedShiftInPicometers}
  \uses{def:physics:phyx-mini-0530:phyxminiproblems-problemphyxmini0530-answerchoice}
  Picometer value printed beside each answer label (1 pm = 10⁻¹² m)
\end{definition}
\begin{proof}
  This is the declared model definition of displayed Shift In Picometers.
\end{proof}
\begin{definition}[recorded Dataset Answer]
  \label{def:physics:phyx-mini-0530:phyxminiproblems-problemphyxmini0530-recordeddatasetanswer}
  \lean{PhyXMiniProblems.ProblemPhyXMini0530.recordedDatasetAnswer}
  \uses{def:physics:phyx-mini-0530:phyxminiproblems-problemphyxmini0530-answerchoice}
  Answer label recorded by the source dataset
\end{definition}
\begin{proof}
  This is the declared model definition of recorded Dataset Answer.
\end{proof}
\begin{definition}[Rounds To Nearest Hundredth Picometer]
  \label{def:physics:phyx-mini-0530:phyxminiproblems-problemphyxmini0530-roundstonearesthundredthpicometer}
  \lean{PhyXMiniProblems.ProblemPhyXMini0530.RoundsToNearestHundredthPicometer}
  Agreement with a picometer value displayed to the nearest 0.01 pm
\end{definition}
\begin{proof}
  This is the declared model definition of Rounds To Nearest Hundredth Picometer.
\end{proof}
\begin{definition}[Matches Answer Choice]
  \label{def:physics:phyx-mini-0530:phyxminiproblems-problemphyxmini0530-matchesanswerchoice}
  \lean{PhyXMiniProblems.ProblemPhyXMini0530.MatchesAnswerChoice}
  \uses{def:physics:phyx-mini-0530:phyxminiproblems-problemphyxmini0530-sequentialcomptonsetup, def:physics:phyx-mini-0530:phyxminiproblems-problemphyxmini0530-wavelengthshiftinpicometers, def:physics:phyx-mini-0530:phyxminiproblems-problemphyxmini0530-answerchoice, def:physics:phyx-mini-0530:phyxminiproblems-problemphyxmini0530-displayedshiftinpicometers, def:physics:phyx-mini-0530:phyxminiproblems-problemphyxmini0530-roundstonearesthundredthpicometer}
  A displayed choice agrees with the modeled net wavelength shift
\end{definition}
\begin{proof}
  This is the declared model definition of Matches Answer Choice.
\end{proof}
% --- Archon declaration topology end ---
% --- Archon physics formalization source end ---
